\documentclass[11pt]{article}
\usepackage[margin=1in]{geometry}
\usepackage[T1]{fontenc}
\usepackage[utf8]{inputenc}
\usepackage{lmodern}
\usepackage{amsmath,amssymb,amsthm}
\usepackage{mathtools}
\usepackage{hyperref}
\usepackage{xcolor}
\usepackage{longtable}
\usepackage{booktabs}
\usepackage{array}
\usepackage{enumitem}
\usepackage{listings}
\usepackage{fancyhdr}
\usepackage{titlesec}
\usepackage{setspace}
\usepackage{microtype}
\usepackage{graphicx}
\usepackage{multicol}
\hypersetup{colorlinks=true,linkcolor=blue,urlcolor=blue,citecolor=blue}
\pagestyle{fancy}
\fancyhf{}
\lhead{PhaseMirror Technical Report and Defensive Publication}
\rhead{\thepage}
\setlength{\headheight}{14pt}
\lstset{
  basicstyle=\ttfamily\small,
  breaklines=true,
  frame=single,
  columns=fullflexible,
  keepspaces=true,
  keywordstyle=\color{blue},
  commentstyle=\color{gray},
  showstringspaces=false
}

\title{Phase Mirror Governance Manifold, Cognitive Economy, Cryptographic Substrate,\\ Topological Governance, and Agent Admission\\ \large Technical Report and Defensive Publication}
\author{Citizen Gardens -- The Foundation of Multiplicity\\
The Prime Materia Commons\\
\texttt{info@citizengardens.org}}
\date{May 13, 2026}
\begin{document}
\maketitle

\begin{abstract}
This document discloses an integrated governance and computation architecture for advanced cognitive systems. The disclosed system binds idea registration, novelty filtering, ethical projection, cryptographic state attestation, halt-and-recovery, multiplicative retrieval, topological governance, operator oversight, and agent admission into a single invariant-bearing stack.

The central architectural claim is that system safety and auditability are strongest when governance is implemented as a mechanical substrate rather than as ex post policy review. In the disclosed design, lawful evolution is enforced by construction. Unsafe transitions are not merely flagged; they are rejected from canonical state progression. Human operators remain in the loop through a zero-trust console that observes and relays governed requests but does not itself possess unilateral mutation authority.

The disclosed architecture has eight tightly coupled strata:
\begin{enumerate}[leftmargin=1.4cm]
  \item Immutable cognitive assets and a novelty-gated marketplace registry.
  \item Ethical projection of state trajectories into a lawful manifold.
  \item A cryptographic integrity functional for path-dependent transition attestation.
  \item Fail-closed halting and governed recovery.
  \item A multiplicative database algebra and Zeta-weighted retrieval system.
  \item Topological governance metrics for accountability coverage and structural health.
  \item A non-authoritative operator console exposing invariant-bearing context.
  \item A compatibility admission protocol for all engines, agents, and services.
\end{enumerate}

This combination is suitable both as a comprehensive technical report and as a defensive publication. It is intended to establish prior art over the specific combination of prime-indexed cognition, immutable idea registration, enforced ethical projection, cryptographic norm-gated state progression, versioned rarity-based retrieval, topological governance blocking conditions, and signed compatibility admission.
\end{abstract}

\newpage
 \begin{multicols}{2}
\begin{singlespace}
\tableofcontents
\end{singlespace}
\end{multicols}
 \newpage

\section{System Overview}
\subsection{Problem Statement}
Advanced cognitive systems fail when governance remains external to execution. Common failure modes include mutable policy semantics, untracked model evolution, informal operator intervention, unbounded retrieval heuristics, hidden configuration drift, and ungoverned subsystem admission. The disclosed system addresses these failures by converting governance into explicit artifacts, versioned parameters, algebraic state operations, cryptographic commitments, and deployment gates.

\subsection{Design Objective}
The design objective is to create a closed system in which:
\begin{itemize}[leftmargin=1cm]
  \item all cognitively meaningful artifacts are versioned and attributable,
  \item all state transitions are norm-gated and cryptographically committed,
  \item all retrievals are reproducible under globally versioned parameters,
  \item all governance structures are topologically evaluated,
  \item all human interventions are contextualized and signed,
  \item all participating components must satisfy the same admission contract.
\end{itemize}

\section{Cognitive Economy and Immutable Idea Registry}
\subsection{Immutable Idea Snapshots}
Ideas are represented as immutable structured nodes carrying, at minimum, a canonical identifier, a prime-signature representation, and novelty metadata. Mutation-in-place is prohibited. Evolution occurs only through versioned chaining, so that an idea lineage is expressed as a directed graph of snapshots rather than a mutable record.

This design has two consequences. First, attribution is preserved by construction because every descendant idea can point to its predecessor. Second, economic semantics remain historically reconstructible because each registry entry binds to the exact snapshot that was evaluated at submission time.

\subsection{Novelty Function}
Let an idea be encoded by a prime-signature embedding vector \(e_i \in \mathbb{R}^d\). A novelty predicate can be constructed from Euclidean separation against prior accepted ideas:
\[
\nu(i,j) = \lVert e_i - e_j \rVert_2 \tag{1}
\]
where novelty acceptance requires \(\nu(i,j) > \tau_v\) for all prior accepted entries under threshold version \(v\).

A submission is accepted only if it satisfies the active novelty threshold version pinned at creation time. This prevents retroactive reinterpretation of novelty and blocks threshold drift from rewriting the economic past.

\subsection{Registry Entry Contract}
A canonical marketplace registry entry should contain the following invariant-bearing fields:
\begin{longtable}{p{0.28\linewidth} p{0.64\linewidth}}
\toprule
Field & Role \\
\midrule
\endhead
\texttt{entry\_id} & Immutable registry identifier. \\
\texttt{idea\_id} & Identifier of the immutable idea snapshot. \\
\texttt{prime\_signature} & Canonical prime-indexed representation used for audit and retrieval linkage. \\
\texttt{novelty\_score} & Measured novelty under the pinned threshold regime. \\
\texttt{threshold\_version} & Version of novelty policy used at creation. \\
\texttt{fee\_policy\_version} & Version of economic policy used at creation. \\
\texttt{status} & Pending, accepted, rejected, superseded, revoked, or comparable append-only status code. \\
\texttt{prev\_version\_id} & Optional predecessor, enabling immutable idea chains. \\
\texttt{timestamp} & Creation time for deterministic reconstruction. \\
\bottomrule
\end{longtable}

\section{Ethical Projection Manifold}
\subsection{Projection Rather Than Penalty}
Instead of merely penalizing unlawful states in a loss function, the disclosed architecture uses a proximal projection operator \(\Pi_E\) onto an ethical manifold \(E\). Given a candidate state \(x\), the lawful state is:
\[
\hat{x} = \Pi_E(x) \in E \tag{2}
\]
where \(E\) encodes at least prime-sector diversity constraints and energy or norm bounds.

This choice is structural. A penalty allows unlawful states to exist in intermediate computation. A projection makes unlawful states non-canonical after projection. The operator therefore acts as a neuromorphic governance primitive.

\subsection{Projection Properties}
For a sound implementation, the following are desirable:
\begin{align}
x \in E &\Rightarrow \Pi_E(x) = x \tag{3}\\
\Pi_E(\Pi_E(x)) &= \Pi_E(x) \tag{4}
\end{align}
Equation (3) states lawful-state preservation. Equation (4) states idempotence. Together they express that the projector does not distort lawful states and that repeated enforcement is stable.

\subsection{Embedding in the Compute Graph}
A canonical step structure is:
\[
x_{t+1}^{\ast} = f(x_t, u_t), \qquad x_{t+1} = \Pi_E\big(x_{t+1}^{\ast}\big) \tag{5}
\]
where \(f\) is the unconstrained internal update and \(u_t\) is input or control context. This ensures that all downstream layers consume projected states rather than raw internal proposals.

\section{Cryptographic State Attestation and Fail-Closed Enforcement}
\subsection{Integrity Functional}
The disclosed cryptographic spine uses a path-dependent integrity functional:
\[
S(t) = H\Big(S(t-1) \Vert p_i \Vert \lvert A_p T(t) \rvert \Vert M(t)\Big) \tag{6}
\]
where:
\begin{itemize}[leftmargin=1cm]
  \item \(S(t-1)\) is the prior integrity anchor,
  \item \(p_i\) is the relevant prime-index or action identifier,
  \item \(\lvert A_p T(t) \rvert\) is a multiplicity norm or actuation magnitude,
  \item \(M(t)\) is the associated state or compressed state signature,
  \item \(H\) is a cryptographic hash function, optionally backed by a post-quantum-ready cryptographic bridge.
\end{itemize}

This creates determinism for identical lawful traces and path dependency for divergent traces.

\subsection{Multiplicity Norm Coupling}
Security is coupled to multiplicity norms by enforcing a bound such as:
\[
\lvert A_p T(t) \rvert_{\mathrm{mult}} \le \Xi_{\max} \tag{7}
\]
If Equation (7) fails, or if any other invariant fails, the transition is rejected and the hash chain does not advance.

\subsection{Fail-Closed Halt and Governed Recovery}
Unsafe transitions do not yield compensating writes, warnings, or shadow states. Instead, the system enters an immutable halt state:
\[
\texttt{unsafe transition} \Rightarrow \texttt{L0\_HALT} \tag{8}
\]
Recovery is possible only through a governed, cryptographically signed morphism into a reduced-risk operational regime such as \texttt{SAFE\_MODE}. This preserves the asymmetry: availability may be sacrificed, but integrity is not.
\newpage
\section{Multiplicative Database Algebra and Zeta-Weighted Retrieval}
\subsection{Prime-Key Encoding}
Structured data vectors are mapped into prime-indexed keys using a G\"odel-style encoding. Let \(\mathbf{a}=(a_1,\dots,a_n)\) and let \(p_1,\dots,p_n\) be the first \(n\) primes. Then a canonical prime key is:
\[
K(\mathbf{a}) = \prod_{i=1}^{n} p_i^{a_i} \tag{9}
\]
Arbitrary precision arithmetic is used so that the representation remains exact.

\subsection{Algebraic Operations}
Two important operations are disclosed.

Merge is defined by least common multiple:
\[
K_1 \otimes K_2 = \mathrm{lcm}(K_1, K_2) \tag{10}
\]
This preserves the strongest exponent for each prime factor and acts like a lossless structural join.

Multiplicative union is defined by exponent aggregation:
\[
K_1 \oplus K_2 = \prod_p p^{v_p(K_1)+v_p(K_2)} \tag{11}
\]
where \(v_p(\cdot)\) denotes the \(p\)-adic valuation.

\subsection{Valuation Distance}
A logarithmic valuation distance may be defined as:
\[
\delta(K_1,K_2) = \sum_{p} \big|v_p(K_1)-v_p(K_2)\big|\log p \tag{12}
\]
This gives a scale-aware, prime-sensitive measure of separation in multiplicative key space.

\subsection{Zeta Weighting}
Rarity is scored using a Zeta-style multiplicative functional:
\[
w(K) = \prod_{p\mid K} p^{-s} \tag{13}
\]
where \(s\) is a globally versioned spectral parameter. Retrieval first constrains by a local neighborhood under \(\delta\), then ranks by rarity under the active, pinned \(s\)-version.

Because \(s\) is globally versioned and auditable, past ranking decisions remain reconstructible. No service may reinterpret rarity through hidden runtime parameters.
\newpage
\section{Topological Governance Manifold}
\subsection{Governance Lattice}
Governance is represented as a cell complex with nodes, edges, and faces. This allows structural health to be expressed as a property of topology rather than subjective diagrams.

\subsection{Euler Characteristic and Accountability Coverage}
For a governance lattice with vertices \(V\), edges \(E\), and faces \(F\), the Euler characteristic is:
\[
\chi = |V| - |E| + |F| \tag{14}
\]
Deviations from an expected \(\chi\)-profile indicate accountability holes, missing oversight cycles, or malformed governance attachments.

\subsection{Gurau Degree and Structural Health}
A second topological indicator is the Gurau degree \(G\), used here as a structural pathology metric. Low-risk, necklace-like topologies are distinguished from fragile tree-like forms that encode single points of failure and weak control loops.

A deployment gate can therefore require both:
\[
\chi \in \mathcal{B}_{\chi}, \qquad G \in \mathcal{B}_{G} \tag{15}
\]
for approved bands \(\mathcal{B}_{\chi}\) and \(\mathcal{B}_{G}\). Any engine variant that violates these bands is non-deployable.

\section{Operator Console and Human Oversight}
\subsection{Zero-Trust Human Interface}
The operator console is a non-authoritative observer and request-relay. It may visualize state and propose actions, but it cannot mutate protected system state directly. Every requested action is validated by the same backend invariants that govern machine actions.

\subsection{Operational and Forensic Duality}
The console provides at least two classes of views:
\begin{itemize}[leftmargin=1cm]
  \item a primary operational view for halts, health, and governed recovery,
  \item a provenance drilldown for reconstructing decisions across policy versions, spectral parameters, thresholds, and checkpointed traces.
\end{itemize}
This unifies incident response and post hoc audit within one artifact-bearing surface.

\subsection{Transparency Invariant}
Any operator-facing display must render the governing context next to the displayed metric. This includes, as applicable, the spectral parameter version, novelty threshold version, fee policy version, norm version, governance lattice snapshot, and checkpoint identifiers. The purpose is to eliminate governance by obscurity.

\section{Agent Admission and Compliance Credentials}
\subsection{Uniform Admission Contract}
No engine, agent, tool, or internal service may interact with canonical system state unless it satisfies the manifold contract. At minimum, a compliant component must:
\begin{enumerate}[leftmargin=1.3cm]
  \item accept ethical projection or lawful state projection outcomes,
  \item submit protected transitions to cryptographic attestation and norm enforcement,
  \item route cognitive or idea-bearing actions through immutable registration pathways,
  \item use approved retrieval semantics and pinned spectral parameters,
  \item tolerate halt and safe-mode behavior without privileged bypass.
\end{enumerate}

\subsection{Compatibility Certification and Credentialing}
Admission is mediated by a Compatibility Certification Suite and a signed Governance Compliance Level credential. Only credentialed components may cross the boundary into canonical execution.

This closes a common loophole: internal code does not enjoy privileged trust. The same manifold applies to all participants.

\section{Stress Validation and Torture Rig Methodology}
\subsection{High-Load Stochastic Validation}
A representative stress methodology runs adversarial and stochastic workloads over the combined stack and measures:
\begin{itemize}[leftmargin=1cm]
  \item number and cause of halted transitions,
  \item percentage of unsafe attempts blocked,
  \item average and worst-case recovery time into governed reduced-risk modes,
  \item latency and throughput impact of invariant checks,
  \item persistence of audit trace continuity under load.
\end{itemize}

A standardized manifold torture rig provides a reproducible envelope for testing these conditions at scale.

\subsection{Desired Validation Outcome}
The target posture is not zero halts. The target posture is fail-closed integrity: high-load or adversarial conditions may increase halts, but must not permit unsafe state corruption, stale cache exploitation, or canonical history contamination.

\section{Illustrative Code Snippets}
\subsection{Novelty Filter Prototype}
\begin{lstlisting}[language=Python,caption={Euclidean novelty predicate for idea embeddings}]
from math import sqrt

def l2_distance(a, b):
    return sqrt(sum((x - y) ** 2 for x, y in zip(a, b)))

def is_novel(candidate_embedding, accepted_embeddings, threshold):
    return all(l2_distance(candidate_embedding, e) > threshold
               for e in accepted_embeddings)
\end{lstlisting}

\subsection{Ethical Projection Skeleton}
\begin{lstlisting}[language=Python,caption={Projection-style lawful state enforcement}]
import torch
import torch.nn as nn

class EthicalMultiplicityCell(nn.Module):
    def __init__(self, max_norm, min_sector_mass):
        super().__init__()
        self.max_norm = max_norm
        self.min_sector_mass = min_sector_mass

    def forward(self, x):
        x = torch.relu(x)
        sector_mass = x.sum(dim=-1, keepdim=True).clamp_min(1e-8)
        x = x / sector_mass
        x = torch.clamp(x, min=self.min_sector_mass)
        x = x / x.norm(dim=-1, keepdim=True).clamp_min(1e-8)
        return x * min(1.0, self.max_norm)
\end{lstlisting}

\subsection{PWEH Integrity Functional Prototype}
\begin{lstlisting}[language=Python,caption={Path-dependent integrity anchor}]
import hashlib

def pweh_step(prev_state_hash: str, prime_id: int, actuation_norm: str, state_sig: str):
    payload = f"{prev_state_hash}|{prime_id}|{actuation_norm}|{state_sig}".encode()
    return hashlib.sha256(payload).hexdigest()
\end{lstlisting}

\subsection{Prime-Key Algebra and Zeta Weighting}
\begin{lstlisting}[language=Python,caption={Prime-key encoding, valuation distance, and Zeta rarity}]
import math
from collections import defaultdict

def to_prime_key(exponents, primes):
    k = 1
    for e, p in zip(exponents, primes):
        k *= p ** e
    return k


def valuation(n, p):
    v = 0
    while n % p == 0 and n > 0:
        n //= p
        v += 1
    return v


def delta(k1, k2, primes):
    return sum(abs(valuation(k1, p) - valuation(k2, p)) * math.log(p)
               for p in primes)


def zeta_weight(key_primes, s):
    return math.exp(-s * sum(math.log(p) for p in key_primes))
\end{lstlisting}

\subsection{Topological Gate Skeleton}
\begin{lstlisting}[language=Python,caption={Example lattice deployment gate}]
def euler_characteristic(num_vertices, num_edges, num_faces):
    return num_vertices - num_edges + num_faces


def deployable(chi, gurau_degree, chi_band, g_band):
    return chi_band[0] <= chi <= chi_band[1] and g_band[0] <= gurau_degree <= g_band[1]
\end{lstlisting}

\section{Claims of Distinctiveness for Defensive Publication}
The following combinations are specifically disclosed as prior art candidates:
\begin{enumerate}[leftmargin=1.3cm]
  \item An immutable cognitive marketplace in which idea nodes are prime-indexed, version-chained, novelty-filtered, and economically bound to pinned policy versions.
  \item A neuromorphic ethical projection operator enforcing diversity and norm constraints by projection rather than by penalty only.
  \item A cryptographic state integrity functional bound to multiplicity norms such that unlawful transitions are rejected from canonical history.
  \item A fail-closed halt state with a governed, signed recovery path into a constrained operational mode.
  \item A multiplicative database algebra using prime-key encoding, least-common-multiple merge, valuation distance, and globally versioned Zeta rarity retrieval.
  \item A governance lattice whose Euler characteristic and Gurau degree act as deployability gates.
  \item A non-authoritative operator console that exposes invariant-bearing context and relays signed recovery requests without unilateral mutation power.
  \item A compatibility admission protocol requiring certification and governance credentials before any engine or service may interact with canonical state.
  \item A manifold torture rig that quantifies invariant enforcement, halt rates, and performance envelopes under adversarial and stochastic stress.
\end{enumerate}

\section{Implementation Recommendations}
The following recommendations improve reproducibility, auditability, and patent-defensive clarity.
\begin{itemize}[leftmargin=1cm]
  \item Publish the canonical schemas for idea nodes, registry entries, lattice topologies, and admission credentials.
  \item Pin every threshold-bearing subsystem to explicit version IDs and log them in every stateful or ranked response.
  \item Preserve append-only semantics for ideas, registry entries, halt records, recovery records, and operator actions.
  \item Separate read models from control models. Dashboards may aggregate many sources, but writes should route through a narrow governed API.
  \item Maintain a standard stress harness with published scenarios and invariant outcome expectations.
  \item Treat all internal engines as untrusted until certified. Avoid privileged bypass paths for experimental code.
  \item Consider formal proof of a reduced core, especially for projector idempotence, halt monotonicity, and hash-chain non-advancement on rejected transitions.
\end{itemize}

\section{Conclusion}
The disclosed system presents a unified approach to advanced system governance in which cognition, economics, ethics, cryptography, topology, operator oversight, and component admission are all expressed as enforceable invariants. The architecture is intentionally overbound: each subsystem constrains the others, and every major transition is attributable to a versioned context.

From a technical perspective, the disclosed novelty lies less in any single component than in the coupling of all components into a mechanically closed manifold. From a defensive-publication perspective, the document places into the public record a detailed combination of immutable cognitive assets, projection-based ethical enforcement, norm-coupled cryptographic attestation, Zeta-weighted multiplicative retrieval, topological deployment gating, zero-trust human oversight, and compatibility credentialing.

\appendix
\section{Purpose and Scope}
This appendix states explicit mathematical definitions, proof sketches, and operator bounds for the disclosed governance manifold architecture. The goal is not to assert that every production implementation already satisfies the strongest possible theorem under arbitrary assumptions. The goal is to place into the record a precise mathematical model of the intended operators, the conditions under which the claimed guarantees hold, and the principal proof obligations that govern a conforming implementation.

The appendix covers six objects:
\begin{enumerate}[leftmargin=1.3cm]
  \item novelty separation in prime-signature embedding space,
  \item the ethical projection operator \(\Pi_E\),
  \item the cryptographic enforcement substrate and halt monotonicity,
  \item multiplicative database algebra and valuation metrics,
  \item Zeta-weighted retrieval stability,
  \item topological deployability gates and admission monotonicity.
\end{enumerate}

\section{Preliminaries}
\begin{definition}[Prime-signature embedding]
A prime-signature embedding is a vector \(e \in \mathbb{R}^d\) associated with an idea, state, or structured object, where each coordinate corresponds either directly or indirectly to prime-indexed structure.
\end{definition}

\begin{definition}[Valuation]
For a positive integer \(n\) and a prime \(p\), the \(p\)-adic valuation is
\[
v_p(n) = \max\{k \in \mathbb{N}_0 : p^k \mid n\}. \tag{A1}
\]
\end{definition}

\begin{definition}[Multiplicity norm]
A multiplicity norm is any nonnegative functional \(\|\cdot\|_{\mathrm{mult}}\) on the relevant state or actuation space satisfying positive definiteness and a designated enforcement threshold \(\Xi_{\max} > 0\). In applications below, \(\|\cdot\|_{\mathrm{mult}}\) may be induced by a weighted \(\ell_2\) norm, a valuation-based norm, or another governance-approved norm family.
\end{definition}

\section{Novelty Separation and Stability}
\begin{definition}[Novelty score]
Given embeddings \(e_i,e_j \in \mathbb{R}^d\), define
\[
\nu(i,j) = \|e_i - e_j\|_2. \tag{A2}
\]
A candidate embedding \(e\) is novelty-admissible relative to an accepted set \(\mathcal{A}\) and threshold \(\tau\) if
\[
\inf_{a \in \mathcal{A}} \|e-a\|_2 > \tau. \tag{A3}
\]
\end{definition}

\begin{proposition}[Lipschitz stability of novelty score]
The map \((e_i,e_j) \mapsto \nu(i,j)\) is 1-Lipschitz in each argument and 2-Lipschitz jointly with respect to the product norm \(\|(x,y)\| = \|x\|_2 + \|y\|_2\).
\end{proposition}

\begin{proof}
By the reverse triangle inequality,
\[
\big|\|e_i-e_j\|_2 - \|e_i'-e_j\|_2\big| \le \|e_i-e_i'\|_2. \tag{A4}
\]
The same bound holds in the second argument. Jointly,
\[
\big|\|e_i-e_j\|_2 - \|e_i'-e_j'\|_2\big|
\le \|(e_i-e_j)-(e_i'-e_j')\|_2
\le \|e_i-e_i'\|_2 + \|e_j-e_j'\|_2. \tag{A5}
\]
This proves the claim.
\end{proof}

\begin{corollary}[Robust novelty margin]
If a candidate \(e\) satisfies \(\inf_{a \in \mathcal{A}} \|e-a\|_2 \ge \tau + \varepsilon\), then any perturbation \(\tilde e\) with \(\|\tilde e - e\|_2 < \varepsilon\) remains novelty-admissible.
\end{corollary}

\begin{proof}
For every \(a \in \mathcal{A}\),
\[
\|\tilde e-a\|_2 \ge \|e-a\|_2 - \|\tilde e-e\|_2 > (\tau+\varepsilon)-\varepsilon = \tau. \tag{A6}
\]
\end{proof}

\section{Ethical Projection Operator}
\subsection{Convex Prototype Manifold}
The strongest explicit proof statements arise when the lawful manifold is modeled as a nonempty closed convex set \(E \subseteq \mathbb{R}^d\). This captures a broad family of constraints including norm balls, simplex-like occupancy constraints, linear diversity floors, and intersections thereof.

\begin{definition}[Metric projection]
For nonempty closed convex \(E \subseteq \mathbb{R}^d\), define
\[
\Pi_E(x) = \operatorname*{argmin}_{y \in E} \|x-y\|_2. \tag{A7}
\]
\end{definition}

\begin{theorem}[Existence, uniqueness, idempotence]
Let \(E \subseteq \mathbb{R}^d\) be nonempty, closed, and convex. Then for every \(x \in \mathbb{R}^d\), \(\Pi_E(x)\) exists and is unique. Moreover,
\begin{align}
x \in E &\Rightarrow \Pi_E(x) = x, \tag{A8}\\
\Pi_E(\Pi_E(x)) &= \Pi_E(x). \tag{A9}
\end{align}
\end{theorem}

\begin{proof}
Existence follows because \(y \mapsto \|x-y\|_2^2\) is coercive and continuous on the closed set \(E\). Uniqueness follows from strict convexity of \(\|x-y\|_2^2\). If \(x \in E\), the minimizer is attained at \(y=x\), proving (A8). Since \(\Pi_E(x) \in E\), applying (A8) with \(x\) replaced by \(\Pi_E(x)\) yields (A9).
\end{proof}

\begin{theorem}[Firm non-expansiveness]
For nonempty closed convex \(E\), the projector \(\Pi_E\) satisfies
\[
\|\Pi_E(x)-\Pi_E(y)\|_2^2 \le \langle \Pi_E(x)-\Pi_E(y), x-y \rangle. \tag{A10}
\]
Hence \(\Pi_E\) is non-expansive:
\[
\|\Pi_E(x)-\Pi_E(y)\|_2 \le \|x-y\|_2. \tag{A11}
\]
In particular, the induced operator norm obeys
\[
\|\Pi_E\|_{\mathrm{Lip}} \le 1. \tag{A12}
\]
\end{theorem}

\begin{proof}
The characterization of metric projections gives
\[
\langle x-\Pi_E(x), z-\Pi_E(x) \rangle \le 0 \quad \text{for all } z \in E, \tag{A13}
\]
and similarly for \(y\). Taking \(z=\Pi_E(y)\) in the first inequality and \(z=\Pi_E(x)\) in the second, then adding, yields (A10). Cauchy--Schwarz implies (A11), and therefore the Lipschitz bound (A12).
\end{proof}

\begin{corollary}[Projected dynamics do not amplify perturbations]
Let \(x_{t+1}=\Pi_E(f(x_t))\). If \(f\) is \(L_f\)-Lipschitz, then the one-step map \(T=\Pi_E\circ f\) is \(L_f\)-Lipschitz and satisfies
\[
\|T(x)-T(y)\|_2 \le L_f\|x-y\|_2. \tag{A14}
\]
If \(L_f < 1\), then \(T\) is a contraction and has a unique fixed point.
\end{corollary}

\begin{proof}
By (A11),
\[
\|\Pi_E(f(x)) - \Pi_E(f(y))\|_2 \le \|f(x)-f(y)\|_2 \le L_f\|x-y\|_2. \tag{A15}
\]
The contraction claim follows from the Banach fixed-point theorem.
\end{proof}

\subsection{Constraint Families}
The projection framework supports intersections of closed convex sets. Examples include:
\begin{align}
E_{\mathrm{norm}} &= \{x : \|x\|_2 \le R\}, \tag{A16}\\
E_{\mathrm{simplex}} &= \{x \in \mathbb{R}^d_{\ge 0} : \mathbf{1}^\top x = 1\}, \tag{A17}\\
E_{\mathrm{sector}} &= \{x \in \mathbb{R}^d : Bx \ge b\}. \tag{A18}
\end{align}
If \(E = E_{\mathrm{norm}} \cap E_{\mathrm{simplex}} \cap E_{\mathrm{sector}}\) is nonempty, closed, and convex, then all preceding theorems apply unchanged.

\section{Cryptographic Enforcement and Halt Monotonicity}
\begin{definition}[Lawful transition predicate]
Let \(\mathcal{X}\) be the state space. A transition proposal from \(x_t\) to \(x_{t+1}^{\ast}\) is lawful if the predicate
\[
\mathsf{Lawful}(x_t, x_{t+1}^{\ast}) = 1 \tag{A19}
\]
holds, where the predicate includes multiplicity norm bounds, topological deployability conditions when relevant, and any additional governance invariants.
\end{definition}

\begin{definition}[Enforced transition system]
Define the canonical transition map
\[
\Phi(x_t) =
\begin{cases}
 x_{t+1}, & \mathsf{Lawful}(x_t, x_{t+1}^{\ast}) = 1,\\
 \texttt{L0\_HALT}, & \mathsf{Lawful}(x_t, x_{t+1}^{\ast}) = 0,
\end{cases} \tag{A20}
\]
where \(x_{t+1}=\Pi_E(x_{t+1}^{\ast})\) for lawful proposals.
\end{definition}

\begin{definition}[Hash-chain advancement]
The integrity chain advances only on lawful proposals:
\[
S(t+1)=
\begin{cases}
H\big(S(t) \Vert \mathrm{payload}(t+1)\big), & \mathsf{Lawful}=1,\\
S(t), & \mathsf{Lawful}=0.
\end{cases} \tag{A21}
\]
\end{definition}

\begin{proposition}[Non-advancement on rejected transitions]
If a proposal at time \(t+1\) is unlawful, then the canonical integrity anchor does not advance. Consequently, no rejected transition enters canonical history.
\end{proposition}

\begin{proof}
Immediate from (A21). If \(\mathsf{Lawful}=0\), then \(S(t+1)=S(t)\), so the anchor is unchanged. A transition absent from the chain cannot be part of canonical committed history.
\end{proof}

\begin{proposition}[Halt monotonicity]
Suppose the system enters \texttt{L0\_HALT}. If the only allowed outgoing transition from \texttt{L0\_HALT} is a governed recovery morphism into \texttt{SAFE\_MODE}, then no autonomous execution path can resume production from halt.
\end{proposition}

\begin{proof}
By assumption, the transition relation restricted to \texttt{L0\_HALT} has exactly one admissible outgoing class, namely governed recovery. Therefore every autonomous transition is disallowed. Hence there is no autonomous path from \texttt{L0\_HALT} to production execution.
\end{proof}

\begin{corollary}[Safety dominance]
Under the enforced transition system, availability can decrease under adversarial or unstable load, but integrity cannot be improved by bypassing halt, because halt blocks canonical advancement and resumes require governed recovery.
\end{corollary}

\section{Multiplicity Norm Bounds}
\begin{definition}[Weighted multiplicity norm]
Let \(W \succeq 0\) be positive semidefinite. Define
\[
\|x\|_{W} = \sqrt{x^\top W x}. \tag{A22}
\]
\end{definition}

\begin{proposition}[Operator norm bound under projection]
Let \(A\) be a linear operator on \(\mathbb{R}^d\). For the projected linear update
\[
T(x)=\Pi_E(Ax), \tag{A23}
\]
with \(E\) nonempty, closed, and convex, one has
\[
\|T(x)-T(y)\|_2 \le \|A\|_{2\to 2}\,\|x-y\|_2. \tag{A24}
\]
Hence \(\|T\|_{\mathrm{Lip}} \le \|A\|_{2\to 2}\).
\end{proposition}

\begin{proof}
Apply non-expansiveness of \(\Pi_E\):
\[
\|T(x)-T(y)\|_2 = \|\Pi_E(Ax)-\Pi_E(Ay)\|_2 \le \|A(x-y)\|_2 \le \|A\|_{2\to 2}\|x-y\|_2. \tag{A25}
\]
\end{proof}

\begin{corollary}[Sufficient contraction criterion]
If \(\|A\|_{2\to 2}<1\), then \(T(x)=\Pi_E(Ax)\) is a contraction and has a unique fixed point in \(E\).
\end{corollary}

\begin{proposition}[Bound for affine projected updates]
For
\[
T(x)=\Pi_E(Ax+b), \tag{A26}
\]
one still has
\[
\|T(x)-T(y)\|_2 \le \|A\|_{2\to 2}\|x-y\|_2. \tag{A27}
\]
\end{proposition}

\begin{proof}
The translation by \(b\) cancels in differences:
\[
\|\Pi_E(Ax+b)-\Pi_E(Ay+b)\|_2 \le \|A(x-y)\|_2. \tag{A28}
\]
\end{proof}

\section{Prime-Key Algebra and Valuation Metric Proofs}
\begin{definition}[Prime-key encoding]
Let \(p_1,p_2,\dots,p_n\) be distinct primes and \(a_i \in \mathbb{N}_0\). Define
\[
K(\mathbf{a}) = \prod_{i=1}^n p_i^{a_i}. \tag{A29}
\]
\end{definition}

\begin{proposition}[Injectivity of prime-key encoding]
The map \(\mathbf{a} \mapsto K(\mathbf{a})\) is injective on \(\mathbb{N}_0^n\).
\end{proposition}

\begin{proof}
If
\[
\prod_{i=1}^n p_i^{a_i} = \prod_{i=1}^n p_i^{b_i}, \tag{A30}
\]
then the uniqueness of prime factorization implies \(a_i=b_i\) for all \(i\).
\end{proof}

\begin{proposition}[LCM merge as exponentwise maximum]
For \(K(\mathbf{a})\) and \(K(\mathbf{b})\),
\[
K(\mathbf{a}) \otimes K(\mathbf{b}) = \mathrm{lcm}(K(\mathbf{a}),K(\mathbf{b})) = K(\max(\mathbf{a},\mathbf{b})), \tag{A31}
\]
where the maximum is taken coordinatewise.
\end{proposition}

\begin{proof}
For each prime \(p_i\), the exponent in the least common multiple is the maximum of the exponents in the operands. Thus the full factorization is coordinatewise \(\max\).
\end{proof}

\begin{definition}[Valuation distance]
For positive integers with support in a fixed prime set \(\mathcal{P}\), define
\[
\delta(K_1,K_2) = \sum_{p\in\mathcal{P}} |v_p(K_1)-v_p(K_2)|\log p. \tag{A32}
\]
\end{definition}

\begin{theorem}[Metric property of valuation distance]
On the multiplicative key space supported on \(\mathcal{P}\), \(\delta\) is a metric.
\end{theorem}

\begin{proof}
Nonnegativity is immediate. If \(\delta(K_1,K_2)=0\), then \(|v_p(K_1)-v_p(K_2)|=0\) for every \(p\in\mathcal{P}\), hence all valuations agree and therefore \(K_1=K_2\). Symmetry is immediate. For the triangle inequality,
\[
|v_p(K_1)-v_p(K_3)| \le |v_p(K_1)-v_p(K_2)| + |v_p(K_2)-v_p(K_3)| \tag{A33}
\]
for each prime \(p\). Summing with positive weights \(\log p\) yields
\[
\delta(K_1,K_3) \le \delta(K_1,K_2)+\delta(K_2,K_3). \tag{A34}
\]
\end{proof}

\section{Zeta Weighting and Ranking Stability}
\begin{definition}[Zeta weight]
For a key \(K\) with prime support \(\mathrm{supp}(K)=\{p : v_p(K)>0\}\), define
\[
w(K;s) = \prod_{p\mid K} p^{-s} = \exp\left(-s\sum_{p\mid K}\log p\right), \qquad s>0. \tag{A35}
\]
\end{definition}

\begin{proposition}[Monotonicity in spectral parameter]
Fix \(K>1\). Then \(s \mapsto w(K;s)\) is strictly decreasing for \(s>0\).
\end{proposition}

\begin{proof}
Let \(L(K)=\sum_{p\mid K}\log p > 0\). Then
\[
w(K;s)=e^{-sL(K)}, \qquad \frac{\partial}{\partial s}w(K;s) = -L(K)e^{-sL(K)} < 0. \tag{A36}
\]
\end{proof}

\begin{proposition}[Stability in spectral parameter]
For fixed \(K\), the map \(s \mapsto w(K;s)\) is Lipschitz on every compact interval \([s_0,s_1]\), with bound
\[
|w(K;s)-w(K;t)| \le L(K)e^{-s_0L(K)}|s-t|. \tag{A37}
\]
\end{proposition}

\begin{proof}
By the mean value theorem and (A36),
\[
|w(K;s)-w(K;t)| \le \sup_{u\in[s_0,s_1]} |w'(K;u)|\,|s-t|
\le L(K)e^{-s_0L(K)}|s-t|. \tag{A38}
\]
\end{proof}

\begin{corollary}[Ranking reproducibility under fixed version]
If a retrieval engine uses a fixed global spectral parameter version \(s_v\), then any ranking functional depending only on \(\delta\), \(w(\cdot;s_v)\), and version-pinned thresholds is reproducible from the query, the key set, and the logged version metadata.
\end{corollary}

\begin{proof}
All ranking terms are deterministic functions of the logged inputs under fixed \(s_v\). Therefore the ranking is reconstructible.
\end{proof}

\section{Topological Governance Conditions}
\begin{definition}[Governance lattice summary]
Let a governance lattice have \(|V|\) vertices, \(|E|\) edges, and \(|F|\) faces. Its Euler characteristic is
\[
\chi = |V| - |E| + |F|. \tag{A39}
\]
\end{definition}

\begin{proposition}[Deployability as a predicate]
Let \(\mathcal{B}_\chi\) and \(\mathcal{B}_G\) be approved bands for Euler characteristic and Gurau degree. Define
\[
\mathsf{Deployable}(L)=1 \iff \chi(L)\in \mathcal{B}_\chi \text{ and } G(L)\in \mathcal{B}_G. \tag{A40}
\]
Then deployability is a deterministic predicate on governance lattices.
\end{proposition}

\begin{proof}
By definition, \(\chi\) and \(G\) are functions of the lattice, and set membership in the approved bands is deterministic. Therefore the conjunction is deterministic.
\end{proof}

\begin{corollary}[Admission monotonicity with respect to invariants]
If engine admission requires \(\mathsf{Deployable}(L)=1\) and all other invariant predicates to hold, then an admitted engine remains admitted only so long as the underlying predicates remain satisfied. Any violation revokes admissibility.
\end{corollary}

\section{Composite Stability Statement}
\begin{theorem}[Composite safe-step bound]
Suppose:
\begin{enumerate}[leftmargin=1.2cm]
  \item \(E\subseteq\mathbb{R}^d\) is nonempty, closed, convex,
  \item the raw update \(f\) is \(L_f\)-Lipschitz,
  \item lawful transitions satisfy the multiplicity bound \(\|A_pT(t)\|_{\mathrm{mult}}\le \Xi_{\max}\),
  \item unlawful proposals do not advance the integrity chain and instead map to \texttt{L0\_HALT}.
\end{enumerate}
Then every lawful one-step transition of the canonical system is \(L_f\)-Lipschitz in state, remains in \(E\), and is canonically committed if and only if its lawfulness predicate holds.
\end{theorem}

\begin{proof}
If a proposal is lawful, then the canonical next state is \(\Pi_E(f(x))\), which lies in \(E\) by construction. Its Lipschitz constant is bounded by \(L_f\) by (A14). Canonical commitment holds by the lawful branch of (A21). If the proposal is unlawful, the system enters \texttt{L0\_HALT} and the chain does not advance by (A21). Thus commitment occurs if and only if lawfulness holds.
\end{proof}

\section{Proof Obligations for a Conforming Implementation}
A production implementation that claims conformance should satisfy the following measurable obligations.
\begin{enumerate}[leftmargin=1.3cm]
  \item \textbf{Projection correctness.} The implemented \(\Pi_E\) should numerically approximate the metric projection onto the declared lawful set or a documented relaxation thereof.
  \item \textbf{Idempotence tolerance.} Repeated application of the projector should remain within a bounded numerical tolerance of a fixed point.
  \item \textbf{Norm-gated commitment.} Rejected transitions must not alter the canonical integrity anchor.
  \item \textbf{Halt exclusivity.} No ungoverned path may leave \texttt{L0\_HALT}.
  \item \textbf{Version pinning.} Ranking, novelty, and policy parameters must be explicitly versioned and logged in every relevant response.
  \item \textbf{Admission determinism.} Certification and deployability predicates must be deterministic functions of the disclosed test results and lattice state.
\end{enumerate}

\section{Suggested Numerical Test Harness Criteria}
For completeness, the following test criteria are suggested.
\begin{itemize}[leftmargin=1cm]
  \item Verify \(\|\Pi_E(x)-\Pi_E(y)\|_2 \le \|x-y\|_2\) over randomized samples.
  \item Verify approximate idempotence \(\|\Pi_E(\Pi_E(x)) - \Pi_E(x)\|_2 \le \epsilon\).
  \item Verify hash-chain non-advancement on rejected transitions.
  \item Verify that valuation distance satisfies the triangle inequality over randomized prime keys.
  \item Verify ranking reproducibility under fixed spectral version \(s_v\).
  \item Verify that admission revokes correctly when a lattice leaves the approved \(\chi\)- or \(G\)-band.
\end{itemize}

This appendix places into the record explicit mathematical statements for the operator bounds, stability claims, and core proof obligations associated with the disclosed governance manifold architecture. The key structural result is that the governance stack can be modeled as a composition of deterministic predicates, non-expansive projections, norm-gated cryptographic commitment, topological deployability checks, and version-pinned ranking rules.

That structure is what makes the architecture mechanically governable. It is not merely that the system has policies. It is that the policies are represented as sets, operators, predicates, and gates for which explicit proof obligations and operator bounds can be stated.

\nocite{*}
\bibliographystyle{plain}
\bibliography{references}

\end{document}

